\documentclass[12pt]{article}

\usepackage[a4paper, total={6in, 8in}]{geometry}
\usepackage{textcomp}

\begin{document}
\setlength{\parindent}{0pt}

\def \t {\textrightarrow}
\def \v {\vspace{1em}}
\def \bi {\begin{itemize}}
\def \ei {\end{itemize}}
\def \s[#1] {\section*{#1}}
\def \ss[#1] {\subsection*{#1}}
\def \sss[#1] {\subsubsection*{#1}}

\s[Italo Calvino]
Gadda è un ingengiere, e anche Calvino non ha formazione classica \\
Calvino è il canone del 900 \t e sperimenta la stretta connessione tra il mondo della scienza e quello della letteratura, della filologia \\
Lui ha il merito di rappresentare il canone del 900 \\
La leggerezza è guardare con quell atteggiamento propositivo gli aspetti + dolenti della nostra esistenza \t anche cio che + drammatico, rimane drammatico ma cambia prospettiva \\
Ha formazione non tradizionale, ed è il personaggio che ci mette davanti a una forma di neo illuminismo \t ovvero una rinnovata fiducia nella luce della ragione \\
L'autore passa attraverso gli studi di matematica e di logica \t e ha studiato specificamente il calcolo combinatorio \t lui studia le possibilita del reale con il calcolo combinatorio \\
Nasce a Santiago (in Cile) \t poi si trasferisce in liguria e studia agraria \t si laurea con una tesi su Conrad, autore di "Out of darkness" (che si collega alla selva oscura dantesca)

\v

Ha fase prima piu neorealista, come in "Sentieri dei nidi di ragno" \t dove pero c e la guerra letta con gli occhi di qualcun altro \\
Forte osservazione della realta è punto di forza \t opera del 47 (i nidi), sul dopoguerra \t le immagini del romanzo rispondo a quell abbassamento tonale e di narrazione tipico dei personaggi presi dalla realta (si riconduce a verga) \\
L immagine del neorealismo viene data dal quadro 900 \t nel neorealismo i personaggi corrispondo alla realta \t questo si vede in quest opera \\
La trilogia
\bi
  \item barone rampante
  \item ?
  \item visconte dimezzato
\ei
Qua c'è lettura allegorica della reatla \t qua c è critica dei valori della societa, visconte dimezzato e barone rampante \t la realta non risponde alle esigenze dell essere umano \\
Questo non è pero il punto focale di calvino \t che è invece "Le lezioni americane", in cui si hanno veramente le sue riflessioni

\v

Nelle lezioni \t si descrive un mondo cambiato \t non c'è rapporto come intende Rousseau tra uomo e natura \t le citta che non sono + espressione di un rapporto tra uomo e natura \\
Questo mondo riflette anche lui una visione che va sempre + scollandosi dall input tipico del partito comunista italiano \t si vede infatti una serie di autori che si sono allineati ai valori del PCI, anche per opposizione all'esperienza fascista \\
Calvino non trova + corrispondenza di valori, con la realta \t legami fragili \t riporta ai legami dell ermetismo che la realta ha interrotto, e quindi Montale con figure incomprensibili \\
"Marcolvaldo" nel 58, e il fallimento della politica intesa come Aristotele concepi il ruolo della politica nella societa \t l animale sociale non ha + i parametri in cui riconoscersi

\ss[Giornata di uno scrutatore]
"Giornata di uno scrutatore" \t aveva avuto un soggiorno a Parigi, e opera esce nel 63 \t questi sono gli anni in cui l uomo sempbra poter raggiungere oggni successo per boom economico \\
Ma dietro boom econ. c'è umanita dimenticata e la solitudine \t e lui la coglie in un fatto politico \t con le elezioni truccate che vedono l affermarsi di un nuovo partito?? \\
PCI si concentra sulle masse, mentre dall altra parte ci sono gli eredi del fascismo \\
Amerigo (scoperta dell america) viene mandato a fare lo scrutatore \t qua la carica umana è di insegnamento agli altri, ed p un mondo in cui dove la mente non puo, non è possibile concepire la mancanza di rispetto \\
I voti confluiscono a favore della democrazia cristiana, sono stati taroccati e acquisiti con un inganno \t gli scrutatori vengono infatti sostituiti da membri del clero ?? \\
Il mondo di questo testo è un mondo rinchiuso, che presente non i lati dolenti della condizione di una minorita, ma la pars costruens \\
Il messaggio si sintetizza nel finale dell opera \t dopo aver mostrato l opersotia di questa citta del dolore, si pone la domanda: \textbf{esiste una citta perfetta? \t si quella di si incontra l imperfezione} \\
\textbf{La citta perfetta è quella che presenta un umanita palpitante (di cui parlera pasolini)} \\
Il personaggio scopre di fatti l america, ovvero un luogo chiuso lontano dall america (amerigo) \t ma una luce cambia il personaggio: i ritratti di queste persone, che vengono chiamate con identificazioni ?? non per nome \\
Ma la loro operosita fa dimenticare la loro unicita \t non esiste + allora la loro unicita, ma cio che hanno da costruire e da mostrare \t cio che non viene visto, ma è \\
Viaggio attraverso dolore, verso l umano \t fatto di sfaccettature che non possono essere enumerate \\
L autore da giudizio crudo su questo mondo e su: 
\bi
  \item da giudizio nei confronti del sistema politico che penalizza liberta
  \item sugli inservienti, che si preoccupano solo di salvare la faccia ma non delle persone
  \item sulla pragmaticita e aletterarietà \t è un saggio dentro un opera letteraria
\ei
Qua siamo nei paraggi di Torino, a Cotolengo \\
L opera è pero anche caratterizzata dalla leggerezza \t per questo ha carattere educativo

\v

"Il castello dei destini incrociati" \t nel 68? \t figure che si districano in un mondo di combinazioni, ma non scientifiche, ma nelle carte dei tarocchi \t le possibilita dei tarocchi \\
Poi nel 72 "Le citta invisibili" \t con carrellata di citta, in tutto 60 \\
"Se una notte di inverno un viaggiatore" del 79 \\
1883, poi "Palomar" = insieme di racconti, su fatti minuti ed episodi piccoli della realta \t ordinati secondo una combinazione che ricorda i processi del calcolo combinatorio \\
L obbiettivo comune è spiegare la realta 

\v

La linea è l intellettuale che si scolla da una formazone prettamente classica \t e ha molteplicita di interessi, spinto dalla curiositas \\
Anche lui rappresenta l intellettuale mediatico \\
"Le formiche argetine" \t parlano di un calvino che non trascura il mondo animale, e lo piega a spiegare una societa come quella del maturo 900 \\

\s[Pasolini]
Animi rozzi, ma autentici \t un umanita dimenticata \t intellettuale che è anche mediatico \t quando si parla di umanita viscerale, e dell intellettuale del 900 \t non si puo dimenticare pasolini \\
È un personaggio demistificato, e ha parabola di vita intensa \t in cui la cultura sua vive metamorfosi continue, e di lui non si puo usare un termine che lo circoscriva (ma una rosa di terminini che provano a descriverlo) \\
Con Pasolini si parla di cinema, di poesia, di scritti, interviste televisive, borgate romane, carcere, e una morte violenta \\
Ci sono testimonianze di pasolini teatrale, regista, mediatico \t c'è anche una trasposizione del decamerone \\
Pasolini guarda la realta e ne urla il dolore \t lo si chiama cantore dell umanita dimenticata, intellettuale dell impegno civile, il demistificatore della nuova societa del profitto e della pubblicita \\
Questo non è il benessere dell anima \t subisce giudizio stroncante sulla sua omosessualita
\end{document}
